% chapters/06/praktikum.tex
\begin{praktikumbox}
Praktikum ini mengikuti Tugas Praktikum 04 pada RPS: membuat prototipe
\istilah{platformer} mini yang bendanya jatuh karena gravitasi, melompat sekali
ketika berada di tanah, dan memantul secara mendatar. Isilah lembar kerja pada
Tabel~\ref{tab:lembar-praktikum-04}.

\begin{langkah}
  \item \textbf{Jalankan program fisika.} Ketik ulang \code{ch06_fisika.py}
        dari Daftar~\ref{lst:ch06-fisika} ke dalam folder \code{code}.
        Jalankan, lalu amati perilaku kedua benda: pemain yang berhenti di
        lantai dan bola yang memantul makin pendek.

  \item \textbf{Uji lompatan.} Tekan tombol lompatan beberapa kali. Pastikan
        pemain hanya melompat ketika menyentuh lantai, dan pastikan menahan
        tombol tidak menghasilkan lompatan berulang.

  \item \textbf{Bandingkan tinggi lompatan dengan hitungan.} Jalankan program
        dengan perintah berikut pada tiga laju bingkai yang berbeda, lalu catat
        tinggi terukur dan tinggi menurut hitungan.
\begin{kodekeluaran}
python ch06_fisika.py --frames 60 --auto-jump 1 --fps 30
python ch06_fisika.py --frames 120 --auto-jump 1 --fps 60
python ch06_fisika.py --frames 240 --auto-jump 1 --fps 120
\end{kodekeluaran}

  \item \textbf{Ubah rasa lompatan.} Jalankan program dengan tiga pasangan
        nilai pada Tabel~\ref{tab:rasa-lompatan}, yaitu
        \code{--gravity 1200 --jump-speed 600},
        \code{--gravity 1800 --jump-speed 735}, dan
        \code{--gravity 2400 --jump-speed 849}. Tingginya sama,
        tetapi rasa geraknya berbeda.

  \item \textbf{Ubah koefisien restitusi bola.} Jalankan program dengan
        \code{--restitution 0.5}, lalu \code{--restitution 0.95}, lalu
        \code{--restitution 1.0}. Amati lama bola memantul pada masing-masing
        nilai, dan catat pengamatan Anda.

  \item \textbf{Tambahkan satu bongkahan untuk dipijak.} Tambahkan satu persegi
        panjang yang berfungsi sebagai lantai kedua pada ketinggian tertentu,
        lalu atur agar pemain dapat mendarat di atasnya. Untuk tahap ini,
        cukup periksa posisi pemain terhadap bongkahan itu secara langsung.

  \item \textbf{Tulis catatan praktikum} yang memuat tabel pengukuran tinggi
        lompatan, hasil percobaan koefisien restitusi, dan kendala yang Anda
        hadapi pada langkah keenam.

  \item \textbf{Simpan pekerjaan} ke repositori Git dengan pesan penyimpanan
        yang menerangkan isinya.
\end{langkah}
\end{praktikumbox}

\begin{table}[htbp]
\centering
\caption{Lembar kerja Praktikum 04}
\label{tab:lembar-praktikum-04}
{\setlength{\tabcolsep}{4pt}\footnotesize%
\begin{tabularx}{\textwidth}{@{}C{1.8cm} C{2.6cm} C{2.6cm} Y@{}}
\toprule
\headrow \textbf{Batas laju} & \textbf{Tinggi terukur} & \textbf{Tinggi menurut hitungan} & \textbf{Selisih} \\
\midrule
30 FPS & & & \\[1.4em]
\addlinespace
60 FPS & & & \\[1.4em]
\addlinespace
120 FPS & & & \\[1.4em]
\addlinespace
\multicolumn{4}{@{}p{\textwidth}@{}}{\footnotesize\itshape
Gravitasi 1800 piksel/detik kuadrat, kecepatan lompatan 700 piksel/detik,
sehingga tinggi menurut hitungan adalah 136,1 piksel.} \\
\addlinespace
\multicolumn{4}{@{}p{\textwidth}@{}}{\footnotesize\itshape
Banyak pantulan bola pada restitusi 0,5: \dotfill \quad
pada 0,95: \dotfill} \\[1.4em]
\bottomrule
\end{tabularx}}
\end{table}

\begin{catatabox}
Langkah keenam sengaja ditinggalkan tanpa petunjuk rinci. Pemeriksaan posisi
terhadap bongkahan secara langsung memang terbatas: pemain hanya dapat mendarat
di atasnya, tetapi tidak dapat menabraknya dari samping. Keterbatasan itulah
yang menjadi titik tolak Bab~VII.
\end{catatabox}
